\section{Mixed moments for Catalan's constant and \texorpdfstring{$\pi^2$}{pi squared}}\label{sec:moments}

\subsection{Keep the two periods in a prescribed ratio}

A rational relation between $1$, $G=\beta(2)$ and $\pi^2$, with a
nonzero coefficient of $G$, would make $G+\lambda\pi^2$ rational for
some rational $\lambda$. This suggests retaining the two periods in a
fixed ratio inside each matrix entry, before taking a determinant.
Since $\pi^2$ is irrational, proving that $G+\lambda\pi^2$ is
irrational for every rational $\lambda$ would establish independence
of $1,G,\pi^2$. The detailed decay calculation below treats the single
test case $\lambda=1$.
The moments of OpenAI~\cite[Section~2]{OpenAI} make that operation
particularly direct. The adaptation below is from \rec{R11}.

Put $c_j=4^{-j}\binom{2j}{j}$ for nonnegative integers $j$, and $c_j=0$
otherwise. For a polynomial $P(t)=\sum P_it^i$, define
\[
 A_j(P)=\sum_iP_i c_{(j-i-1)/2},\qquad
 B_j(P)=\sum_iP_i c_{(i-j)/2}.
\]
Thus $A_j(P)$ is the coefficient of $t^j$ in
$tP(t)/\sqrt{1-t^2}$ at zero, whereas $B_j(P)$ is the coefficient of
$t^j$ in the polynomial part at infinity
\[
 \Omega(P)=\left[\frac{P(t)}{\sqrt{1-t^{-2}}}\right]_+.
\]
The square roots are expanded with constant term one. Define the moments
\begin{align*}
 M(P,j)&=\int_{-1}^1\int_0^1
  \frac{|t|P(t)s^j}{\sqrt{1-t^2}(1-ts)}\dd s\dd t,\\
 Z(D,j)&=\int_0^1\int_0^1\frac{D(t)s^j}{1-ts}\dd s\dd t
\end{align*}
and, for $Q(s)=\sum Q_js^j$, write
$M(P,Q)=\sum Q_jM(P,j)$ and $Z(D,Q)=\sum Q_jZ(D,j)$.
Their period coefficients are
\begin{equation}\label{eq:moment-periods}
 M(P,j)=M^0(P,j)+4B_j(P)G+\tfrac32A_j(P)\zeta(2),\qquad
 Z(D,j)=Z^0(D,j)+D_j\zeta(2),
\end{equation}
where the superscript $0$ denotes the rational part. These are the
moment identities in \cite[(17)]{OpenAI}.

Suppose $D_0$ has coefficient $A_j(P)$ in every column being used,
and $E(P)$ has coefficient $B_j(P)$ there. Replace $D_0$ by
\[
 D_\lambda=D_0-16\lambda E(P).
\]
Then $M(P,j)-\tfrac32Z(D_\lambda,j)$ has the form
\begin{equation}\label{eq:mixed-entry}
 \text{rational}+4B_j(P)(G+\lambda\pi^2).
\end{equation}
The factor $16$ supplies $24\lambda B_j\zeta(2)=4\lambda B_j\pi^2$.
For example, $P=t^2$ gives $\Omega(P)=t^2+\tfrac12$. In columns
$0\le j<4$, take $D_0=t^3$ and $E=t^2+\tfrac12$. The new polynomial
is $t^3-16\lambda t^2-8\lambda$, and the coefficients of
$G+\lambda\pi^2$ are $2,0,4,0$.

\subsection{Rows, columns and the canonical matrix}

We now assemble $n$ matched rows and $n$ column polynomials into a square
matrix. Each entry will be rational plus a rational multiple of
$\theta_\lambda=G+\lambda\pi^2$. Replacing $\theta_\lambda$ by a variable
$X$ makes its determinant a rational polynomial of degree at most $n$.
Our goal is to make the value at $\theta_\lambda$ small enough to survive
clearing the polynomial's denominators.

The columns use powers $s^i$ in a band $b\le i<L$. Matching the period
coefficients in this band fixes the target; the $b$ lower powers remain
available for arithmetic corrections. Zeros of orders $h$ at $t=1$ and
$q$ at $s=1$ improve the real estimates. A further parameter $g$ counts
pairs of rows that will supply fixed rational rows.

The cutoffs are related by $L=n+b+q$, $C=L+g$ and $H=C+h$.
Here $L$ is the first unused column degree, $C-1$ is the common power
of $t$ in the row formulas, and $H-1$ bounds the row degrees. We use
the following fixed proportions throughout
Sections~\ref{sec:moments}--\ref{sec:corrections}:
\begin{equation}\label{eq:moment-scale}
 \begin{aligned}
 n&=48N,& b&=7N,& q=g&=4N,& h&=2N,\\
 L&=59N,& C&=63N,& H&=65N.&&
 \end{aligned}
\end{equation}
Only $N$ tends to infinity. This choice leaves room for the corrections
while balancing the denominator cost against endpoint vanishing; the
comparison of the two rates will be made in Section~\ref{sec:canonical}.
The rows and columns inherited from~\cite{OpenAI} are
\begin{align}
 P_r(t)&=(1-t)^h t^{C-1}T_{|r-g|}(1/t),\nonumber\\
 D_r(t)&=\sgn(r-g)(1-t)^h t^{C-1}U_{|r-g|-1}(1/t),\label{eq:moment-rows}\\
 Q_j(s)&=s^{b+j}(1-s)^q \qquad(0\le r,j<n),\nonumber
\end{align}
with $D_g=0$. Here $T_d,U_d$ are Chebyshev polynomials. The contact
identity of \cite[(5)]{OpenAI} is
\[
 tP_r(t)/\sqrt{1-t^2}-D_r(t)=O(t^L).
\]
Since the column polynomials use only the raw powers $s^i$ with
$b\le i<L$, this supplies exactly the coefficient matching required
in \eqref{eq:mixed-entry}.

The factor $(1-t)^h$ is useful analytically and should survive the
period matching. Simply truncating $\Omega(P_r)$ to the required band
would lose this factor. Instead, replace each monomial by a polynomial
with the same low coefficients and the required zero at $t=1$:
\begin{equation}\label{eq:completion-basis}
 F_i(t)=(1-t)^h\left[\frac{t^i}{(1-t)^h}\right]_{<L}
 \quad(0\le i<L),\qquad
 E_r=\sum_{i=b}^{L-1}B_i(P_r)F_i.
\end{equation}
Here $[\ ]_{<L}$ means Taylor truncation. Each $F_i$ agrees with $t^i$
below degree $L$ and is divisible by $(1-t)^h$. Thus $E_r$ has the
right coefficients and retains the endpoint zero. We call this choice
of $E_r$ the \emph{canonical completion}.

The actual entries are now explicit:
\begin{align*}
 \mathcal M_{rj}(\lambda)
 &=M(P_r,Q_j)-\tfrac32Z(D_r,Q_j)+24\lambda Z(E_r,Q_j)\\
 &=a_{rj}(\lambda)+b_{rj}\theta_\lambda,
\end{align*}
where $a_{rj}(\lambda)$ is the rational part and
\[
 b_{rj}=4\sum_{i=b}^{L-1}[s^i]Q_j(s)\,B_i(P_r).
\]
Define $M_{N,\lambda}(X)=(a_{rj}(\lambda)+b_{rj}X)_{r,j<n}$.
Its determinant is the canonical determinant polynomial. The endpoint
zeros are already built into this matrix; the integer adjustments
needed for the denominator estimate are introduced next.

\subsection{High completions and low arithmetic corrections}

There are two different freedoms in completing these coefficients.
Changing the terms above degree $L-1$ under the same degree and endpoint
conditions has no effect on the determinant. Changing terms below
degree $b$ can change it, and will be needed for the arithmetic estimate.
The following proposition explains the first assertion.

\begin{proposition}\label{prop:completion-invariance}
Let $E_1,E_2$ be linear functions of $P$, each divisible by
$t^b(1-t)^h$, of degree at most $H-1$, and agreeing with $\Omega(P)$
in degrees $b,\ldots,L-1$. Their mixed determinants are identical,
as polynomials in the target and in $\lambda$.
\end{proposition}
\begin{proof}
The difference belongs to
$V=t^L(1-t)^h\Q[t]_{<g}$. Pair rows $r$ and $2g-r$ for $0\le r<g$.
Their $P$ polynomials agree and their $D$ polynomials have opposite
signs. Subtracting the rows leaves the rational row
$3Z(D_{2g-r},Q_j)$, independent of the target and of the completion.
The polynomials $D_{2g-r}$ form a basis of $V$: after removing
$(1-t)^h$, their distinct lowest degrees are $L,\ldots,C-1$.
Every change of completion therefore adds a combination of these fixed
rows to the remaining rows. All these row operations preserve the
determinant.
\end{proof}

For the second freedom, set $\lambda=1$ and $\theta=G+\pi^2$.
Abbreviate $M_{N,1}(X)$ to $M_N(X)=A_N+XB_N$. Replace $E_r$ by
\[
 E_r+\sum_{i=0}^{b-1}U_{ri}F_i,\qquad U\in\Z^{n\times b}.
\]
The added moments are rational, because the low $F_i$ have zero
coefficients in the column band. With
\begin{equation}\label{eq:corrected-pencil}
 W_{ij}=24Z(F_i,Q_j),\qquad
 P_{N,U}(X)=\det(M_N(X)+UW),
\end{equation}
we obtain a polynomial of degree at most $n$ with rational coefficients.
The entries of $U$ will vary independently; Proposition~\ref{prop:completion-invariance}
does not make these low corrections disappear.
Thus two changes have different jobs. The canonical completion restores
the endpoint zero while matching the target. The additional matrix $UW$
changes only the rational parts of the entries, allowing us to impose
prime divisibility conditions. Its real effect on the determinant is
the unresolved part of the argument.

\subsection{One arithmetic bound for an entire congruence class}

It is useful to impose the denominator savings first and leave a whole
congruence class of corrections available for the real estimate.
The following result is the uniform conclusion of \rec{R01}.

\begin{theorem}\label{thm:arithmetic-cost}
There are integer matrices $U_N^0$ and positive integers $a_N$ such that,
with
\[
 m_N=2^C\prod_{2\sqrt H<p\le H,\ p\text{ odd}}p^2,
\]
every $U=U_N^0+m_N Z$, $Z\in\Z^{n\times b}$, has entries divisible
by $2^C$ and satisfies $a_NP_{N,U}(X)\in\Z[X]$. Moreover
\begin{align}
 \log a_N&\le (A_{\rm arith}+o(1))n^2,\label{eq:arithmetic-rate}\\
 A_{\rm arith}&=\frac{33079}{18432}+\frac{505}{4608}\log2
 =1.87061400307785858968\ldots,\nonumber\\
 \log m_N&\le\left(\frac{63\log2+130}{48}+o(1)\right)n.\nonumber
\end{align}
The same $a_N$ works for every integer $Z$, regardless of its size.
One may choose $|U^0_{ri}|\le m_N/2$.
\end{theorem}

Here are the valuation arguments behind the constant. They also identify
which parts are inherited from the moment construction.
At $2$, the smoothing identities of \cite[(22)--(27)]{OpenAI} split
each raw column into a term with valuation $C-O(\log H)$, a multiple
of its polynomial-part coefficient column, and a $D$ term. Denote that
multiple by $e_1(X)$: it is the source's fixed degree-one polynomial
over the $2$-adic numbers, whose coefficient valuations are bounded below
independently of $N$. Only this valuation bound is needed here.

To keep the column indices visible, let
$\boldsymbol\Omega=(B_i(P_r))_{0\le r<n,\ b\le i<L}$ be the
$n$ by $(L-b)$ coefficient matrix of the polynomial parts. Multiplying
it on the right combines these coefficient columns. For the mixed
construction, its original
contribution and the canonical completion combine as
\[
 \boldsymbol\Omega\bigl(e_1(X)I+24Z_F\bigr),\qquad
 (Z_F)_{ij}=Z^0(F_i,j),\quad b\le i,j<L.
\]
All entries of $Z_F$ have valuation at least $-2\lfloor\log_2H\rfloor$.
The low correction, being divisible by $2^C$, joins the smoothed term.
This grouping retains the coefficient-band index when the determinant
is expanded.

Write the coefficients of $P_r/(1-t)^h$ and $D_r/(1-t)^h$ in descending
degrees as vectors $p_u,d_u$. The Chebyshev recurrence gives their
entry valuations at least $u-1$. The scalar coefficient multiplying
$p_u$ in band position $i$ of $\boldsymbol\Omega$ has valuation at least
$u+i-H+1$. If a determinant term uses $v$ columns from the $\boldsymbol\Omega$
group and $w$ from the $D$ group, alternation forces distinct band
indices and distinct indices within each coefficient family. Summing
their least possible values gives
\[
 v_2(P_{N,U})\ge
 \min_{v,w}\left\{-(H-b)v+\tfrac32v^2+\tfrac12w^2
                   +C(n-v-w)\right\}-O(n\log H).
\]
The minimum is over integers $v,w\ge0$ with $v+w\le n$ and has
$w=n-v$. With $(H-b-n)/n=5/24$, it is
$-505n^2/4608-O(n\log n)$. The two alternations are responsible for
the quadratic saving.

For odd $p>2\sqrt H$, the moment reductions in \cite[Section~3]{OpenAI}
have at most two denominator layers. Multiplying the low moment matrix
by $p^2$ exposes a unit minor. If that minor has rank $r$, solving for
the low coefficients modulo $p^2$ kills $r$ columns and leaves cost at
most $2(n-r)$. The monomial part of $F_i$ supplies the minor; its
endpoint tail has rank at most $h$. This gives ranks $b-h$ for
$11N<p\le H/2$ and $b$ for $H/2<p\le145N/4$.

For $145N/4<p\le52N$, use the $b$ raw columns
$j=p,\ldots,p+b-1$. Their low block reduces to $24I_b$ after
multiplication by $p^2$. Put $K=L-p$, $J=H-p$. After killing this
block, $K+J-2b$ columns retain double poles, and the remaining
simple-pole residues span at most $b$ dimensions. The cost is therefore
$2(K+J-2b)+b=227N-4p$. For $p>52N$ no correction is required;
the same residue count is
\[
 2\bigl(K_++(J-b)_+\bigr)
 +(\min(b,J)-\max(0,K))_+.
\]
These calculations give the following piecewise cost $Nd(p/N)$.
The entry $96N$ includes the range where no useful rank saving is used.
For $2\sqrt H<p\le11N$, choose the zero correction modulo $p^2$.
This incurs the full cost $2n=96N$ and will also supply the divisibility
needed for nonvanishing.
\begin{center}
\begin{tabular}{c r@{\qquad}c r}
\toprule
$x=p/N$&$d(x)$&$x=p/N$&$d(x)$\\
\midrule
$(0,11]$&96&$(145/4,52]$&$227-4x$\\
$(11,65/2]$&86&$(52,58]$&$182-3x$\\
$(65/2,145/4]$&82&$(58,59]$&$124-2x$\\
&&$(59,65]$&$65-x$\\
\bottomrule
\end{tabular}
\end{center}
The integral of $d$ is $33079/8$. The prime number theorem, followed by
division by $n^2=2304N^2$, gives $33079/18432$. Small odd primes cost
$O(n\sqrt H\log H)=o(n^2)$; primes above $H$ cost nothing. The Chinese
remainder theorem combines the local choices with divisibility by
$2^C$. This proves Theorem~\ref{thm:arithmetic-cost}, using the stated
moment reductions. Their exact finite checks are indexed in
Appendix~\ref{app:audit}.

\subsection{Nonvanishing under a rationality assumption}

The arithmetic congruences also permit an independent nonvanishing
argument. This matters because a sequence of zero integers gives no
contradiction, however good its upper bound.
At a prime scale $N=p>260$, we have $H=65p$ and
$2\sqrt{65p}<p\le11p$. Thus $p$ lies in the range where the prescribed
correction is zero modulo $p^2$. Since $p^2\mid m_N$, every matrix
$U=U_N^0+m_NZ$ in the arithmetic class satisfies $p^2\mid U$.
This is the hypothesis used in the following proposition.

\begin{proposition}\label{prop:moment-nonvanishing}
Fix rational $\lambda$. Outside a finite set of primes, at scale $N=p$
the mixed determinant
$P_{N,\lambda,U}(X)=\det(M_{N,\lambda}(X)+UW)$ with $p^2\mid U$ satisfies
\[
 p^{96p}P_{p,\lambda,U}(X)\in\Z_{(p)}[X],
 \qquad p^{96p}P_{p,\lambda,U}(X)\bmod p
       \text{ is a nonzero constant}.
\]
Consequently, for every fixed rational $x$, these determinants are
nonzero at $x$ for all sufficiently large primes. For $\lambda=1$,
every correction in Theorem~\ref{thm:arithmetic-cost} is allowed.
\end{proposition}

The proof in \rec{R04} reduces the unbounded sizes to three polynomial
seed determinants. To specify them, take the scale $N=1$, allow the
extra row $r=48$, and form the $49$ by $48$ rational matrix
\[
 T_\eta(r,k)=\sum_{j=0}^4(-1)^j\binom4j
 \left\{M^0(P_r,7+k+j)-\tfrac32Z^0(D_r,7+k+j)
                  +24\eta Z^0(E_r,7+k+j)\right\}.
\]
Let $T_{\eta,0}$ use rows $0,\ldots,47$ and $T_{\eta,1}$ rows
$1,\ldots,48$, and set
$F_\sigma(Y)=\det(T_{Y,0}+\sigma T_{Y,1})$ for $\sigma=0,1,-1$.
The paired rows remove four powers of $Y$, so each degree is at most
$44$. Exact finite-field calculations give
\begin{center}
\begin{tabular}{rrrr}
\toprule
$\sigma$&Prime&Degree / leading coefficient&Number of roots\\
\midrule
0&103&44 / 82&0\\
1&107&44 / 57&0\\
$-1$&103&44 / 5&0\\
\bottomrule
\end{tabular}
\end{center}
This excludes every rational root. If its denominator is prime to the
listed prime, reduction gives a finite-field root. If the denominator
is divisible by that prime, homogenization to degree $44$ leaves the
nonzero leading term. The degree bound is needed for this second case.

To see why these are the correct seeds, put
$\epsilon=(-1)^{(p-1)/2}$ and $K_p(t)=(1-t^2)^{(p-1)/2}$. The
Frobenius identity
\[
 (1-t^{-2})^{-1/2}
 =\epsilon t^{1-p}K_p(t)(1-t^{-2p})^{-1/2}\pmod p
\]
transfers the coefficient matching and its endpoint completion from
scale $p$ to scale one. The source moment digit reductions then give
a block matrix
\[
 A_*^{\mathsf T}\otimes T_{\epsilon\lambda,0}
 +B_*^{\mathsf T}\otimes T_{\epsilon\lambda,1},
 \qquad B_*=A_*\Pi,
\]
where $A_*$ is invertible and $\Pi e_0=0$, $\Pi e_i=e_{p-i}$ for
$1\le i<p$. Thus its determinant, up to a unit and a sign, is
\[
 F_0(\epsilon\lambda)
 F_1(\epsilon\lambda)^{(p-1)/2}
 F_{-1}(\epsilon\lambda)^{(p-1)/2}.
\]
The target coefficient and the correction vanish in this reduction:
the former is integral before scaling, and $p^2\mid U$ while
$v_p(W_{ij})\ge-2$. Excluding the finitely many numerator and denominator
primes of the six specialized seed values proves the proposition.

Under the hypothesis $\theta=a/d\in\Q$, we now have nonzero integers
\[
 d^n a_NP_{N,U}(\theta)
\]
along large prime scales, for every admissible correction. The remaining
task is analytic: choose those corrections so that their determinants
decay faster than $\exp(-A_{\rm arith}n^2)$.
